% ============================================================
%  ECON306: The Economics of Health and Education
%  W01D1 — Course welcome; what is health economics;
%          the six special characteristics; case for govt involvement
%  Date: Wednesday 15 July 2026
%  University of Otago — Semester 2, 2026
%  Lecturer: Austin Henderson
% ============================================================

\documentclass[aspectratio=169, 12pt]{beamer}

% ============================================================
%  ECON306: Shared Preamble — The Economics of Health and Education
%  University of Otago — Semester 2, 2026
%
%  SHARED BY ALL ECON306 DAY DECKS.
%  Input this file immediately after \documentclass in each deck:
%    % ============================================================
%  ECON306: Shared Preamble — The Economics of Health and Education
%  University of Otago — Semester 2, 2026
%
%  SHARED BY ALL ECON306 DAY DECKS.
%  Input this file immediately after \documentclass in each deck:
%    % ============================================================
%  ECON306: Shared Preamble — The Economics of Health and Education
%  University of Otago — Semester 2, 2026
%
%  SHARED BY ALL ECON306 DAY DECKS.
%  Input this file immediately after \documentclass in each deck:
%    \input{../Preambles/econ306_preamble.tex}
%  Do NOT include \documentclass or per-deck metadata here.
% ============================================================

% ─── Theme & Colours ────────────────────────────────────────
\usetheme{default}
\usecolortheme{default}
\usefonttheme{professionalfonts}
\definecolor{OtagoBlue}{HTML}{003057}
\definecolor{OtagoGold}{HTML}{A8923A}
\definecolor{TextGrey}{HTML}{3A3A3A}
\definecolor{AccentBlue}{HTML}{2B6CB0}
\definecolor{LightBg}{HTML}{F7F9FC}

\setbeamercolor{frametitle}{fg=OtagoBlue, bg=white}
\setbeamercolor{title}{fg=OtagoBlue}
\setbeamercolor{subtitle}{fg=OtagoGold}
\setbeamercolor{author}{fg=TextGrey}
\setbeamercolor{date}{fg=TextGrey}
\setbeamercolor{institute}{fg=TextGrey}
\setbeamercolor{normal text}{fg=TextGrey, bg=white}
\setbeamercolor{item}{fg=OtagoBlue}
\setbeamercolor{subitem}{fg=AccentBlue}
\setbeamercolor{block title}{bg=OtagoBlue!12, fg=OtagoBlue}
\setbeamercolor{block body}{bg=OtagoBlue!5, fg=TextGrey}
\setbeamercolor{alertblock title}{bg=OtagoGold!20, fg=OtagoBlue}
\setbeamercolor{alertblock body}{bg=OtagoGold!8, fg=TextGrey}
\setbeamercolor{alerted text}{fg=OtagoGold!80!black}
\setbeamercolor{structure}{fg=OtagoBlue}

\setbeamertemplate{frametitle}{%
	\vspace{0.35em}%
	\textbf{\insertframetitle}\par%
	\vspace{0.05em}%
	\textcolor{OtagoGold}{\rule{\textwidth}{0.6pt}}}

\setbeamertemplate{navigation symbols}{}
\setbeamertemplate{footline}{%
	\leavevmode\hbox{%
		\begin{beamercolorbox}[wd=.40\paperwidth,ht=2.5ex,dp=1.5ex,leftskip=.3cm]{author in head/foot}%
			\textcolor{TextGrey}{\footnotesize ECON306 \textbullet\ University of Otago}%
		\end{beamercolorbox}%
		\begin{beamercolorbox}[wd=.47\paperwidth,ht=2.5ex,dp=1.5ex,center]{title in head/foot}%
			\textcolor{TextGrey}{\footnotesize \insertshorttitle}%
		\end{beamercolorbox}%
		\begin{beamercolorbox}[wd=.13\paperwidth,ht=2.5ex,dp=1.5ex,rightskip=.3cm,right]{date in head/foot}%
			\textcolor{TextGrey}{\footnotesize \insertframenumber\,/\,\inserttotalframenumber}%
	\end{beamercolorbox}}\vskip0pt}

\setbeamertemplate{itemize item}{\textcolor{OtagoGold}{\small$\blacktriangleright$}}
\setbeamertemplate{itemize subitem}{\textcolor{AccentBlue}{\small$\triangleright$}}
\setbeamertemplate{enumerate item}{\textcolor{OtagoBlue}{\small\bfseries\insertenumlabel.}}

\setbeamertemplate{title page}{%
	\begin{tikzpicture}[remember picture,overlay]
		\fill[OtagoBlue] (current page.north west)
		rectangle ([yshift=-0.38\paperheight]current page.north east);
	\end{tikzpicture}
	\vspace{0.06\paperheight}
	\begin{beamercolorbox}[wd=\textwidth, sep=0pt]{title}
		{\color{white}\Large\textbf{\inserttitle}}\\[0.4em]
		{\color{OtagoGold}\large\insertsubtitle}
	\end{beamercolorbox}
	\vspace{1.4em}
	\begin{beamercolorbox}[wd=\textwidth, sep=0pt]{author}
		{\color{TextGrey}\normalsize\insertauthor}\\[0.2em]
		{\color{TextGrey}\small\insertinstitute}\\[0.2em]
		{\color{TextGrey}\small\insertdate}
\end{beamercolorbox}}

% ─── Packages ──────────────────────────────────────────────
\usepackage{tikz}
\usepackage{booktabs}
\usepackage{array}
\usepackage{graphicx}
\usepackage{hyperref}
\usepackage{amsmath}
\usepackage{amssymb}
\usepackage{colortbl}
\usepackage{multirow}
\usepackage{xcolor}
\hypersetup{colorlinks=true, linkcolor=AccentBlue, urlcolor=AccentBlue,
	citecolor=AccentBlue, pdfauthor={Austin Henderson},
	pdftitle={ECON306 S2 2026}}

% ─── Images folder ─────────────────────────────────────────
%  Place all slide images in the 'images/' subfolder next to each deck.
%  Usage: \includegraphics[width=0.85\textwidth]{figure_name}
\graphicspath{{images/}}

% ─── Reusable macros ───────────────────────────────────────
% Recap slide at start of each lecture
\newcommand{\recapslide}[1]{%
	\begin{frame}{Recap: What Did We Cover Last Time?}
		\begin{block}{Key points from last lecture}
			\begin{itemize}#1\end{itemize}
		\end{block}
\end{frame}}

% Plan slide at start of each lecture
\newcommand{\planslide}[2]{%
	\begin{frame}{Today's Plan — #1}
		\begin{columns}
			\column{0.6\textwidth}
			\begin{itemize}#2\end{itemize}
			\column{0.38\textwidth}
			\begin{alertblock}{Reading}
				\footnotesize See Brightspace (Aoroa) for required readings before this lecture.
			\end{alertblock}
		\end{columns}
\end{frame}}

% Summary slide at end of each lecture
\newcommand{\summaryside}[1]{%
	\begin{frame}{Lecture Summary}
		\begin{exampleblock}{Key takeaways}
			\begin{itemize}#1\end{itemize}
		\end{exampleblock}
		\vspace{0.5em}
		{\footnotesize\textcolor{OtagoGold}{\textbf{Brightspace (Aoroa):}} slides, readings, and any announcements posted there.}
\end{frame}}

% NZ spotlight box
\newcommand{\nzbox}[2]{%
	\begin{alertblock}{\includegraphics[height=0.8em]{nz_flag_icon}\hspace{0.3em} NZ Spotlight: #1}
		#2
\end{alertblock}}

% Tutorial callout
\newcommand{\tutorialbox}[1]{%
	\begin{block}{\textbf{Tutorial This Week}}
		#1
\end{block}}

% ─── Section dividers ──────────────────────────────────────
\AtBeginSection[]{
	\begin{frame}[plain]
		\begin{tikzpicture}[remember picture,overlay]
			\fill[OtagoBlue!95] (current page.north west)
			rectangle (current page.south east);
			\node[anchor=center, text=white, font=\Large\bfseries,
			text width=0.9\paperwidth, align=center]
			at (current page.center) {\insertsectionhead};
			\node[anchor=south, text=OtagoGold, font=\normalsize,
			text width=0.9\paperwidth, align=center]
			at ([yshift=2.5cm]current page.south) {ECON306 — University of Otago};
		\end{tikzpicture}
\end{frame}}

%  Do NOT include \documentclass or per-deck metadata here.
% ============================================================

% ─── Theme & Colours ────────────────────────────────────────
\usetheme{default}
\usecolortheme{default}
\usefonttheme{professionalfonts}
\definecolor{OtagoBlue}{HTML}{003057}
\definecolor{OtagoGold}{HTML}{A8923A}
\definecolor{TextGrey}{HTML}{3A3A3A}
\definecolor{AccentBlue}{HTML}{2B6CB0}
\definecolor{LightBg}{HTML}{F7F9FC}

\setbeamercolor{frametitle}{fg=OtagoBlue, bg=white}
\setbeamercolor{title}{fg=OtagoBlue}
\setbeamercolor{subtitle}{fg=OtagoGold}
\setbeamercolor{author}{fg=TextGrey}
\setbeamercolor{date}{fg=TextGrey}
\setbeamercolor{institute}{fg=TextGrey}
\setbeamercolor{normal text}{fg=TextGrey, bg=white}
\setbeamercolor{item}{fg=OtagoBlue}
\setbeamercolor{subitem}{fg=AccentBlue}
\setbeamercolor{block title}{bg=OtagoBlue!12, fg=OtagoBlue}
\setbeamercolor{block body}{bg=OtagoBlue!5, fg=TextGrey}
\setbeamercolor{alertblock title}{bg=OtagoGold!20, fg=OtagoBlue}
\setbeamercolor{alertblock body}{bg=OtagoGold!8, fg=TextGrey}
\setbeamercolor{alerted text}{fg=OtagoGold!80!black}
\setbeamercolor{structure}{fg=OtagoBlue}

\setbeamertemplate{frametitle}{%
	\vspace{0.35em}%
	\textbf{\insertframetitle}\par%
	\vspace{0.05em}%
	\textcolor{OtagoGold}{\rule{\textwidth}{0.6pt}}}

\setbeamertemplate{navigation symbols}{}
\setbeamertemplate{footline}{%
	\leavevmode\hbox{%
		\begin{beamercolorbox}[wd=.40\paperwidth,ht=2.5ex,dp=1.5ex,leftskip=.3cm]{author in head/foot}%
			\textcolor{TextGrey}{\footnotesize ECON306 \textbullet\ University of Otago}%
		\end{beamercolorbox}%
		\begin{beamercolorbox}[wd=.47\paperwidth,ht=2.5ex,dp=1.5ex,center]{title in head/foot}%
			\textcolor{TextGrey}{\footnotesize \insertshorttitle}%
		\end{beamercolorbox}%
		\begin{beamercolorbox}[wd=.13\paperwidth,ht=2.5ex,dp=1.5ex,rightskip=.3cm,right]{date in head/foot}%
			\textcolor{TextGrey}{\footnotesize \insertframenumber\,/\,\inserttotalframenumber}%
	\end{beamercolorbox}}\vskip0pt}

\setbeamertemplate{itemize item}{\textcolor{OtagoGold}{\small$\blacktriangleright$}}
\setbeamertemplate{itemize subitem}{\textcolor{AccentBlue}{\small$\triangleright$}}
\setbeamertemplate{enumerate item}{\textcolor{OtagoBlue}{\small\bfseries\insertenumlabel.}}

\setbeamertemplate{title page}{%
	\begin{tikzpicture}[remember picture,overlay]
		\fill[OtagoBlue] (current page.north west)
		rectangle ([yshift=-0.38\paperheight]current page.north east);
	\end{tikzpicture}
	\vspace{0.06\paperheight}
	\begin{beamercolorbox}[wd=\textwidth, sep=0pt]{title}
		{\color{white}\Large\textbf{\inserttitle}}\\[0.4em]
		{\color{OtagoGold}\large\insertsubtitle}
	\end{beamercolorbox}
	\vspace{1.4em}
	\begin{beamercolorbox}[wd=\textwidth, sep=0pt]{author}
		{\color{TextGrey}\normalsize\insertauthor}\\[0.2em]
		{\color{TextGrey}\small\insertinstitute}\\[0.2em]
		{\color{TextGrey}\small\insertdate}
\end{beamercolorbox}}

% ─── Packages ──────────────────────────────────────────────
\usepackage{tikz}
\usepackage{booktabs}
\usepackage{array}
\usepackage{graphicx}
\usepackage{hyperref}
\usepackage{amsmath}
\usepackage{amssymb}
\usepackage{colortbl}
\usepackage{multirow}
\usepackage{xcolor}
\hypersetup{colorlinks=true, linkcolor=AccentBlue, urlcolor=AccentBlue,
	citecolor=AccentBlue, pdfauthor={Austin Henderson},
	pdftitle={ECON306 S2 2026}}

% ─── Images folder ─────────────────────────────────────────
%  Place all slide images in the 'images/' subfolder next to each deck.
%  Usage: \includegraphics[width=0.85\textwidth]{figure_name}
\graphicspath{{images/}}

% ─── Reusable macros ───────────────────────────────────────
% Recap slide at start of each lecture
\newcommand{\recapslide}[1]{%
	\begin{frame}{Recap: What Did We Cover Last Time?}
		\begin{block}{Key points from last lecture}
			\begin{itemize}#1\end{itemize}
		\end{block}
\end{frame}}

% Plan slide at start of each lecture
\newcommand{\planslide}[2]{%
	\begin{frame}{Today's Plan — #1}
		\begin{columns}
			\column{0.6\textwidth}
			\begin{itemize}#2\end{itemize}
			\column{0.38\textwidth}
			\begin{alertblock}{Reading}
				\footnotesize See Brightspace (Aoroa) for required readings before this lecture.
			\end{alertblock}
		\end{columns}
\end{frame}}

% Summary slide at end of each lecture
\newcommand{\summaryside}[1]{%
	\begin{frame}{Lecture Summary}
		\begin{exampleblock}{Key takeaways}
			\begin{itemize}#1\end{itemize}
		\end{exampleblock}
		\vspace{0.5em}
		{\footnotesize\textcolor{OtagoGold}{\textbf{Brightspace (Aoroa):}} slides, readings, and any announcements posted there.}
\end{frame}}

% NZ spotlight box
\newcommand{\nzbox}[2]{%
	\begin{alertblock}{\includegraphics[height=0.8em]{nz_flag_icon}\hspace{0.3em} NZ Spotlight: #1}
		#2
\end{alertblock}}

% Tutorial callout
\newcommand{\tutorialbox}[1]{%
	\begin{block}{\textbf{Tutorial This Week}}
		#1
\end{block}}

% ─── Section dividers ──────────────────────────────────────
\AtBeginSection[]{
	\begin{frame}[plain]
		\begin{tikzpicture}[remember picture,overlay]
			\fill[OtagoBlue!95] (current page.north west)
			rectangle (current page.south east);
			\node[anchor=center, text=white, font=\Large\bfseries,
			text width=0.9\paperwidth, align=center]
			at (current page.center) {\insertsectionhead};
			\node[anchor=south, text=OtagoGold, font=\normalsize,
			text width=0.9\paperwidth, align=center]
			at ([yshift=2.5cm]current page.south) {ECON306 — University of Otago};
		\end{tikzpicture}
\end{frame}}

%  Do NOT include \documentclass or per-deck metadata here.
% ============================================================

% ─── Theme & Colours ────────────────────────────────────────
\usetheme{default}
\usecolortheme{default}
\usefonttheme{professionalfonts}
\definecolor{OtagoBlue}{HTML}{003057}
\definecolor{OtagoGold}{HTML}{A8923A}
\definecolor{TextGrey}{HTML}{3A3A3A}
\definecolor{AccentBlue}{HTML}{2B6CB0}
\definecolor{LightBg}{HTML}{F7F9FC}

\setbeamercolor{frametitle}{fg=OtagoBlue, bg=white}
\setbeamercolor{title}{fg=OtagoBlue}
\setbeamercolor{subtitle}{fg=OtagoGold}
\setbeamercolor{author}{fg=TextGrey}
\setbeamercolor{date}{fg=TextGrey}
\setbeamercolor{institute}{fg=TextGrey}
\setbeamercolor{normal text}{fg=TextGrey, bg=white}
\setbeamercolor{item}{fg=OtagoBlue}
\setbeamercolor{subitem}{fg=AccentBlue}
\setbeamercolor{block title}{bg=OtagoBlue!12, fg=OtagoBlue}
\setbeamercolor{block body}{bg=OtagoBlue!5, fg=TextGrey}
\setbeamercolor{alertblock title}{bg=OtagoGold!20, fg=OtagoBlue}
\setbeamercolor{alertblock body}{bg=OtagoGold!8, fg=TextGrey}
\setbeamercolor{alerted text}{fg=OtagoGold!80!black}
\setbeamercolor{structure}{fg=OtagoBlue}

\setbeamertemplate{frametitle}{%
	\vspace{0.35em}%
	\textbf{\insertframetitle}\par%
	\vspace{0.05em}%
	\textcolor{OtagoGold}{\rule{\textwidth}{0.6pt}}}

\setbeamertemplate{navigation symbols}{}
\setbeamertemplate{footline}{%
	\leavevmode\hbox{%
		\begin{beamercolorbox}[wd=.40\paperwidth,ht=2.5ex,dp=1.5ex,leftskip=.3cm]{author in head/foot}%
			\textcolor{TextGrey}{\footnotesize ECON306 \textbullet\ University of Otago}%
		\end{beamercolorbox}%
		\begin{beamercolorbox}[wd=.47\paperwidth,ht=2.5ex,dp=1.5ex,center]{title in head/foot}%
			\textcolor{TextGrey}{\footnotesize \insertshorttitle}%
		\end{beamercolorbox}%
		\begin{beamercolorbox}[wd=.13\paperwidth,ht=2.5ex,dp=1.5ex,rightskip=.3cm,right]{date in head/foot}%
			\textcolor{TextGrey}{\footnotesize \insertframenumber\,/\,\inserttotalframenumber}%
	\end{beamercolorbox}}\vskip0pt}

\setbeamertemplate{itemize item}{\textcolor{OtagoGold}{\small$\blacktriangleright$}}
\setbeamertemplate{itemize subitem}{\textcolor{AccentBlue}{\small$\triangleright$}}
\setbeamertemplate{enumerate item}{\textcolor{OtagoBlue}{\small\bfseries\insertenumlabel.}}

\setbeamertemplate{title page}{%
	\begin{tikzpicture}[remember picture,overlay]
		\fill[OtagoBlue] (current page.north west)
		rectangle ([yshift=-0.38\paperheight]current page.north east);
	\end{tikzpicture}
	\vspace{0.06\paperheight}
	\begin{beamercolorbox}[wd=\textwidth, sep=0pt]{title}
		{\color{white}\Large\textbf{\inserttitle}}\\[0.4em]
		{\color{OtagoGold}\large\insertsubtitle}
	\end{beamercolorbox}
	\vspace{1.4em}
	\begin{beamercolorbox}[wd=\textwidth, sep=0pt]{author}
		{\color{TextGrey}\normalsize\insertauthor}\\[0.2em]
		{\color{TextGrey}\small\insertinstitute}\\[0.2em]
		{\color{TextGrey}\small\insertdate}
\end{beamercolorbox}}

% ─── Packages ──────────────────────────────────────────────
\usepackage{tikz}
\usepackage{booktabs}
\usepackage{array}
\usepackage{graphicx}
\usepackage{hyperref}
\usepackage{amsmath}
\usepackage{amssymb}
\usepackage{colortbl}
\usepackage{multirow}
\usepackage{xcolor}
\hypersetup{colorlinks=true, linkcolor=AccentBlue, urlcolor=AccentBlue,
	citecolor=AccentBlue, pdfauthor={Austin Henderson},
	pdftitle={ECON306 S2 2026}}

% ─── Images folder ─────────────────────────────────────────
%  Place all slide images in the 'images/' subfolder next to each deck.
%  Usage: \includegraphics[width=0.85\textwidth]{figure_name}
\graphicspath{{images/}}

% ─── Reusable macros ───────────────────────────────────────
% Recap slide at start of each lecture
\newcommand{\recapslide}[1]{%
	\begin{frame}{Recap: What Did We Cover Last Time?}
		\begin{block}{Key points from last lecture}
			\begin{itemize}#1\end{itemize}
		\end{block}
\end{frame}}

% Plan slide at start of each lecture
\newcommand{\planslide}[2]{%
	\begin{frame}{Today's Plan — #1}
		\begin{columns}
			\column{0.6\textwidth}
			\begin{itemize}#2\end{itemize}
			\column{0.38\textwidth}
			\begin{alertblock}{Reading}
				\footnotesize See Brightspace (Aoroa) for required readings before this lecture.
			\end{alertblock}
		\end{columns}
\end{frame}}

% Summary slide at end of each lecture
\newcommand{\summaryside}[1]{%
	\begin{frame}{Lecture Summary}
		\begin{exampleblock}{Key takeaways}
			\begin{itemize}#1\end{itemize}
		\end{exampleblock}
		\vspace{0.5em}
		{\footnotesize\textcolor{OtagoGold}{\textbf{Brightspace (Aoroa):}} slides, readings, and any announcements posted there.}
\end{frame}}

% NZ spotlight box
\newcommand{\nzbox}[2]{%
	\begin{alertblock}{\includegraphics[height=0.8em]{nz_flag_icon}\hspace{0.3em} NZ Spotlight: #1}
		#2
\end{alertblock}}

% Tutorial callout
\newcommand{\tutorialbox}[1]{%
	\begin{block}{\textbf{Tutorial This Week}}
		#1
\end{block}}

% ─── Section dividers ──────────────────────────────────────
\AtBeginSection[]{
	\begin{frame}[plain]
		\begin{tikzpicture}[remember picture,overlay]
			\fill[OtagoBlue!95] (current page.north west)
			rectangle (current page.south east);
			\node[anchor=center, text=white, font=\Large\bfseries,
			text width=0.9\paperwidth, align=center]
			at (current page.center) {\insertsectionhead};
			\node[anchor=south, text=OtagoGold, font=\normalsize,
			text width=0.9\paperwidth, align=center]
			at ([yshift=2.5cm]current page.south) {ECON306 — University of Otago};
		\end{tikzpicture}
\end{frame}}


\title[ECON306: W01D1 — Course Introduction]{ECON306\\The Economics of Health and Education}
\subtitle{W01D1: Course Welcome \& The Six Special Characteristics}
\author{Austin Henderson}
\institute{Department of Economics\\University of Otago}
\date{Wednesday 15 July 2026}

\begin{document}

\begin{frame}[plain]
  \titlepage
\end{frame}

% ── Course Admin Frames ────────────────────────────────────
\begin{frame}{Course Information}
  \begin{columns}
    \column{0.55\textwidth}
      \textbf{Lectures} (3 per week)\\[0.3em]
      \begin{tabular}{ll}
        Wednesday & 8:00--8:50am \\
        Thursday  & 5:00--5:50pm \\
        Friday    & 1:00--1:50pm \\
      \end{tabular}\\[0.8em]
      \textbf{Tutorials} (every other week, 6 total)\\
      \begin{tabular}{ll}
        Monday & 10:00--10:50am (55U213) \emph{or} \\
        Friday & 10:00--10:50am (T204) \\
      \end{tabular}\\[0.4em]
      {\footnotesize Starting Week 2 --- six tutorials total}
    \column{0.42\textwidth}
      \begin{block}{Contact}
        \textbf{Austin Henderson}\\
        Room 603, Commerce Building\\
        \href{mailto:austin.henderson@otago.ac.nz}{austin.henderson@otago.ac.nz}\\[0.4em]
        \emph{Office hours:} after class\\
        or by appointment
      \end{block}
      \begin{alertblock}{Online resources}
        \textbf{Brightspace (Aoroa)}\\
        {\footnotesize Slides, readings, grades, announcements}
      \end{alertblock}
  \end{columns}
\end{frame}

\begin{frame}{A Bit About Me}
  \begin{itemize}
    \item From the west coast of the US (Los Angeles, San Diego, Seattle)
    \item Lots of subject matter changes: religious studies, then neuroeconomics, then
          behavioral economics, now health. Life and careers are seldom straightforward.
    \item Pivoted to health economics at the start of the COVID-19 pandemic ---
          I saw an opportunity to put research to use immediately (and it's where there were jobs)
    \item Chief health economist on a large federal grant studying COVID-19 among
          American Indian and Alaska Native (AI/AN) peoples
    \item Moved to Dunedin last year, great timing
  \end{itemize}
\end{frame}

\begin{frame}{Course Overview: Three Parts}
  \begin{columns}
    \column{0.33\textwidth}
      \begin{block}{Part I (Weeks 1--7)}
        \textbf{The Economics of\\Health \& Health Care}\\[0.4em]
        {\small Introduction, demand, government, ageing, insurance, hospitals, NZ system}
      \end{block}
    \column{0.33\textwidth}
      \begin{block}{Part II (Weeks 8--9)}
        \textbf{Health Economic\\Evaluation}\\[0.4em]
        {\small Valuing life, discounting, QALYs, equity, MCDA}
      \end{block}
    \column{0.33\textwidth}
      \begin{block}{Part III (Weeks 10--11)}
        \textbf{The Economics\\of Education}\\[0.4em]
        {\small Government role, human capital, signalling, NZ student loans}
      \end{block}
  \end{columns}
  \vspace{0.6em}
  \begin{center}
    \textbf{Together, health + education $\approx$ 13\% of NZ GDP, heavy government
    involvement, and from an economics perspective, lots of cross-over}
  \end{center}
\end{frame}

\begin{frame}{Assessment}
  \begin{center}
    \begin{tabular}{@{}lrp{6.5cm}@{}}
      \toprule
      \textbf{Component} & \textbf{Weight} & \textbf{Details} \\
      \midrule
      Reading Quizzes & 10\% & Weekly on Brightspace (Aoroa); due \textbf{Wednesday} of
                               each teaching week. 30\% late penalty thereafter. \\[0.3em]
      In-class Test   & 15\% & Covers Part I (Weeks 1--7).
                               \textbf{Friday 28 Aug 2026} (Week 7). \\[0.3em]
      Assignment      & 15\% & PHARMAC-style funding brief (Part II focus).
                               Distributed Week 8.
                               \textbf{Due Mon 5 Oct 2026, 9:00am.} \\[0.3em]
      Final Exam      & 55\% & 2-hour exam; covers all material ---
                               lectures, tutorials, and required readings \\
      \bottomrule
    \end{tabular}
  \end{center}
\end{frame}

\begin{frame}{Assessment — Notes}
  \begin{alertblock}{Plussage applies}
    Grades computed under (a) default weights above and (b) in-class test weighted to 0\%
    with final exam correspondingly higher. You receive the higher of the two grades.
    Ask Austin if unsure.
  \end{alertblock}
  \vspace{0.6em}
  {\small Tests and assignments returned in class or from Economics reception
  (6th floor, OBS), 11am--12pm or 2--3pm.}
\end{frame}

\begin{frame}{Grading Scale}
  \begin{center}
    \begin{tabular}{@{}ccccccc@{}}
      \toprule
      \textbf{A+} & \textbf{A} & \textbf{A-} &
      \textbf{B+} & \textbf{B} & \textbf{B-} & \\
      90--100 & 85--89 & 80--84 & 75--79 & 70--74 & 65--69 & \\
      \midrule
      \textbf{C+} & \textbf{C} & \textbf{C-} &
      \textbf{D} & \textbf{E} & & \\
      60--64 & 55--59 & 50--54 & 40--49 & $<$40 & & \\
      \bottomrule
    \end{tabular}
  \end{center}
\end{frame}

\begin{frame}{Tutorials --- Six Sessions, Every Other Week}
  \small
  Tutorials begin in \textbf{Week 2}; two slots per tutorial week --- attend one.
  Please \emph{prepare} before attending: \textbf{debate and discussion}, not passive note-taking.

  \vspace{0.3em}
  \begin{center}
  \begin{tabular}{@{}clll@{}}
    \toprule
    \textbf{T} & \textbf{Week} & \textbf{Date(s)} & \textbf{Theme} \\
    \midrule
    1 & Week 2  & Mon 20 Jul / Fri 24 Jul  & Wagstaff comparative statics \\
    2 & Week 4  & Mon 3 Aug / Fri 7 Aug    & Healthcare for the elderly debate \\
    3 & Week 6  & Mon 17 Aug / Fri 21 Aug  & ``Bob's Party'' + NZ reform \\
    4 & Week 8  & Mon 7 Sep / Fri 11 Sep   & Test debrief + allocation intuitions \\
    5 & Week 10 & Mon 21 Sep / Fri 25 Sep  & PHARMAC CUA case studies \\
    6 & Week 12 & Mon 5 Oct / Fri 9 Oct    & Life advice --- MCDM with 1000minds \\
    \bottomrule
  \end{tabular}
  \end{center}
\end{frame}

\begin{frame}{Required Readings and Quizzes}
  \begin{itemize}
    \item \textbf{Required Reading}: do it \emph{before} the relevant lecture ---
          posted on \textbf{Brightspace (Aoroa)} under eReserve
    \item \textbf{Reading Quizzes}: weekly quiz on Brightspace based on required readings.
          Due \textbf{Wednesday} of each teaching week. 30\% late penalty thereafter.
    \item \textbf{Further Reading} (optional): extends your knowledge, intended for those
          interested in further study or a career in health economics
  \end{itemize}
  \vspace{0.4em}
  \begin{block}{A note on jargon}
    Health economics has many terms with specialised meanings (e.g.\ ``need'' vs.\ ``want'',
    ``efficiency'' vs.\ ``equity''). Learning the distinctions is often the most important
    part of learning health economics. A full glossary is posted on Aoroa.
  \end{block}
\end{frame}

\begin{frame}{Workload Expectations}
  \begin{block}{18-point course $\Rightarrow$ approximately 12 hours per week}
    \begin{itemize}
      \item 3 hours: lectures
      \item 1 hour: tutorial (every other week, averaged out)
      \item \textbf{8--9 hours: your own reading, study, and assessment work}
      \item Check \textbf{Brightspace (Aoroa)} and your student email regularly
      \item Attendance at all classes is expected and strongly recommended
    \end{itemize}
  \end{block}
\end{frame}

% ── Lecture 1 Content: What is Health Economics? ───────────

\planslide{What is Health Economics? The Six Special Characteristics}{
  \item What is health economics? Definitions and scope
  \item Is health care an \emph{economic good}?
  \item The \textbf{six special characteristics} of health care
  \item Implications for funding and organisation
}

\begin{frame}{What is ``Health Economics''?}
  \begin{columns}
    \column{0.52\textwidth}
      \begin{block}{Three definitions}
        \footnotesize
        \begin{itemize}
          \item \textit{Kobelt (1996):} ``the application of theories, tools and concepts
                of economics to the topics of health and health care''
          \item \textit{NICE (UK):} ``the study or analysis of the cost of using and
                distributing health care resources''
          \item \textit{Wikipedia:} ``concerned with efficiency, effectiveness, value and
                behavior in the production and consumption of health and health care''
        \end{itemize}
      \end{block}
    \column{0.44\textwidth}
      \begin{alertblock}{Economics, in brief}
        The study of how individuals and societies use \textbf{scarce resources}.\\[0.4em]
        Faced with scarcity, how do we make the \textbf{best} choices?
      \end{alertblock}
      \vspace{0.4em}
      \begin{exampleblock}{Why it matters for NZ}
        Health + education $\approx$ \textbf{13\% of NZ's GDP}\\
        Health alone: $\approx$\textbf{\$30.6B} per year in 2024/25 ($\sim$8\% GDP)
      \end{exampleblock}
  \end{columns}
\end{frame}

\begin{frame}{Four Fundamental Questions}
  \begin{enumerate}
    \item Is health care an \textbf{economic good}?\\
          {\small (Is health care a fundamental human right?)}
    \vspace{0.4em}
    \item How does health care \textbf{differ} from other economic goods?\\
          {\small (Is health care ``special''?)}
    \vspace{0.4em}
    \item Can we study health and health care using \textbf{economic analysis}?\\
          {\small (Supply and demand?)}
    \vspace{0.4em}
    \item What are the \textbf{implications} of these special characteristics for how health
          care is funded and organised?
  \end{enumerate}
  \vspace{0.5em}
  \begin{alertblock}{A running theme throughout ECON306}
    These questions will recur in every week of the course.
  \end{alertblock}
\end{frame}

\begin{frame}{Special Characteristic 1: Uncertainty}
  \begin{itemize}
    \item Classic economics assumes \emph{perfect information} --- but health care violates this
    \item \textbf{Demand-side uncertainty}: How much HC will you need? When? What type?
          You do not know in advance. This is fundamentally unlike demand for laptops or t-shirts.
    \item \textbf{Supply-side uncertainty}: Doctors make decisions under uncertainty --- what
          is the right diagnosis? The right treatment? Is there a divergence between what is
          best for a patient vs.\ what is best for a population with scarce resources?
    \item \textbf{Insurance} (public or private) is the \emph{institutional response} to this
          uncertainty --- it pools risk across many people
    \item But insurance itself creates new problems: \textit{moral hazard} and
          \textit{adverse selection} (Week 5)
  \end{itemize}
\end{frame}

\begin{frame}{Special Characteristic 2: Derived Demand}
  \begin{itemize}
    \item People don't want health \emph{care} for its own sake --- they want \textbf{health}
    \item HC is therefore a \textit{derived demand}: it is demanded only because it (sometimes)
          produces health
    \item Policy implication: if other inputs (e.g.\ diet, exercise, clean air) produce health
          more cost-effectively than HC, \emph{that is where resources should go}
    \item Demand signals may be weaker when demand is derived --- consumer behaviour is more
          passive
  \end{itemize}
  \vspace{0.3em}
  \begin{block}{Health Production Function (preview)}
    $$H = f(\text{HC},\; \text{nutrition},\; \text{housing},\; \text{education},\; \text{environment},\; \ldots)$$
    Changing relative prices of any input affects the optimal mix.
  \end{block}
\end{frame}

\begin{frame}{Special Characteristic 3: Asymmetric Information}
  \begin{itemize}
    \item \textbf{Principal-agent problem}: providers know more than patients ---
          are their incentives perfectly aligned?
    \item \textbf{Supplier-Induced Demand (SID) hypothesis}: doctors may induce demand for
          their own services (profit motive, NZ uniquely allows direct-to-consumer pharma ads)
    \item On the insurance side: patients who know they are high-risk seek comprehensive cover
          --- \textit{adverse selection}
    \item Societal responses: licensing, credentialing, standardised treatment protocols,
          transparency laws
  \end{itemize}
  \begin{exampleblock}{NZ note}
    NZ and the US are the only developed economies that allow \textbf{direct-to-consumer (DTC)
    pharmaceutical advertising}. This has implications for SID that we examine in Week 3.
  \end{exampleblock}
\end{frame}

\begin{frame}{Special Characteristics 4--5}
  \begin{block}{4. Catastrophic and irreversible consequences}
    \begin{itemize}
      \item Health outcomes can be life-threatening at the extremes of the distribution
      \item Many HC decisions are \textbf{irreversible} --- health is very path-dependent
      \item Risk of catastrophe may matter more to individuals than expected-value averages
    \end{itemize}
  \end{block}
  \vspace{0.4em}
  \begin{block}{5. Externalities}
    \begin{itemize}
      \item \textit{Physical}: immunisation, communicable diseases, food safety, sanitation
      \item \textit{Psychic (caring)}: we care about the health of others (altruism, social norms)
    \end{itemize}
  \end{block}
\end{frame}

\begin{frame}{Special Characteristic 6: Price Discrimination}
  \begin{block}{6. Price discrimination}
    \begin{itemize}
      \item First-degree (perfect), second-degree (bulk/bundles), third-degree (by group)
      \item Insurers use third-degree extensively; regulation governs this in most countries
    \end{itemize}
  \end{block}
  \vspace{0.6em}
  \begin{alertblock}{Summary: Why these six matter}
    Together, these characteristics mean health care markets \textbf{fail} in ways that
    standard markets do not --- setting up the case for government involvement (next slide).
  \end{alertblock}
\end{frame}

\begin{frame}{Implications: The Case for Government Involvement}
  \begin{itemize}
    \item These six characteristics together create \textbf{market failures} in health care ---
          unregulated markets will not produce efficient or equitable outcomes
    \item Three forms of government involvement follow:
    \begin{enumerate}
      \item \textbf{Funding} (e.g.\ tax-financed public insurance, subsidies)
      \item \textbf{Provision} (e.g.\ public hospitals, publicly employed doctors)
      \item \textbf{Regulation} (e.g.\ licensing, pharmaceutical controls, drug pricing)
    \end{enumerate}
    \item In NZ: government funds $\approx$77\% of total health spending
  \end{itemize}
  \vspace{0.3em}
  \begin{alertblock}{But government can also fail}
    Market failure $\Rightarrow$ government intervention \emph{may} improve things, but can also
    make things worse. We will critically examine both throughout ECON306.
  \end{alertblock}
\end{frame}

\begin{frame}{Course Objectives}
  \begin{enumerate}
    \item Apply microeconomic tools (from ECON201/271) to health and education
    \item Understand the history and contemporary landscape of health and education policy
          in NZ, especially its challenges
    \item Introduce \emph{new} microeconomic tools and concepts as needed (e.g.\ QALYs, MCDA,
          Value of Statistical Life)
    \item Equip you to \textbf{understand and critique} economic evaluations of health care
          interventions, and to apply these skills to \emph{any} project appraisal
    \item Develop analytical and decision-making skills, including quantitative proficiencies
  \end{enumerate}
\end{frame}

\summaryside{
  \item Health economics applies economic theory to the production and consumption of health and health care
  \item Health care has \textbf{six special characteristics}: uncertainty, derived demand, asymmetric information, catastrophic consequences, externalities, and price discrimination
  \item These characteristics create \textbf{market failures} that justify government involvement in financing, provision, and regulation
  \item In NZ, health care $\approx$8\% of GDP; assessment: quizzes (10\%), test (15\%), assignment (15\%), exam (55\%)
}

\end{document}
